% Generated draft: requires mathematical and source review.
\documentclass[11pt]{article}
\usepackage[margin=25mm]{geometry}
\usepackage{fontspec}
\setmainfont{Latin Modern Roman}
\usepackage{amsmath,amssymb,mathtools,unicode-math}
\setmathfont{Latin Modern Math}
\usepackage{enumitem,needspace,xurl,hyperref}
\setlength{\parindent}{0pt}
\setlength{\parskip}{5pt}
\setlength{\emergencystretch}{3em}
\newcommand{\N}{\mathbb{N}}
\newcommand{\Z}{\mathbb{Z}}
\newcommand{\Q}{\mathbb{Q}}
\newcommand{\R}{\mathbb{R}}
\newcommand{\C}{\mathbb{C}}
\newcommand{\F}{\mathbb{F}}
\newcommand{\A}{\mathbb{A}}
\newcommand{\PP}{\mathbb{P}}
\newcommand{\E}{\mathbb{E}}
\newcommand{\eps}{\varepsilon}
\newcommand{\dd}{\,\mathrm{d}}
\begin{document}
\section*{UnsolvedMath ID: 11000107}
\textbf{Source problem code:} catalogue entry
\subsection{Definitions and mathematical statement}
Let \(S\) be the source's surface and \(G\le\operatorname{Mod}_S\) any finitely generated subgroup. The notation \(\Phi_f(G)\) denotes an associated group in the source, but its definition is not supplied in this extract. The research summary's tentative ``fixed point / finitely presented kernel'' interpretation is speculation, not a definition, and is not adopted.

For any group \(A\), set \(\gamma_1(A)=A\) and \(\gamma_{j+1}(A)=[A,\gamma_j(A)]\). It is nilpotent when \(\gamma_{c+1}(A)=\{1\}\) for some finite \(c\). The conjecture on PDF page 88 states
\[
 \forall G\le\operatorname{Mod}_S\text{ finitely generated}\quad
 \exists c=c(G)<\infty:\quad
 \gamma_{c+1}(\Phi_f(G))=\{1\}.
\]
The nilpotency class may depend on \(G\); no uniform bound is claimed. The quantifier is over every finitely generated subgroup, not only finite-index or finitely presented subgroups. The source emphasizes limited understanding of finite-index subgroups as motivation, and calls this an old problem, referring to [I4], Section 10.10. A source-complete conjecture requires the definition of \(\Phi_f\) and the surface hypotheses from that reference.
\subsection{Short English statement}
For every finitely generated mapping-class subgroup, the source-defined group \(\Phi_f(G)\) is conjectured nilpotent, with no uniform class bound stated. Its definition is missing, so the summary's tentative interpretation is not used to invent the target.
\subsection{Sources}
\begin{itemize}
\item Local UnsolvedMath registry, permanent ID \(11000107\), exact source code \texttt{catalogue entry}, unnumbered Conjecture, PDF page 88. No individual proposer or proposed year is supplied.
\item Farb, edited volume \emph{Problems on Mapping Class Groups and Related Topics}: \url{https://www.math.uchicago.edu/~farb/papers/mcgbook.pdf}. The volume contains 329 explicitly typeset Problem, Conjecture and Question headings, with page-specific disambiguation used here.
\item Internal reference [I4], Section 10.10, for further discussion. The bibliography entry and definition of \(\Phi_f\) are absent; no author or title is guessed.
\item Extraction pdf-font-structure, accessed 2026-07-29. Source evidence NEEDS\_REVIEW presents an open problem and requires a current primary-literature audit. Registry status open and classification OPEN-TRIAGE; the summary explicitly cannot verify the target and offers only a tentative interpretation of its symbol.
\item Dataset attribution: Ulam AI / UnsolvedMath Contributors, CC BY 4.0. Source rights NEEDS\_REVIEW: publicly accessible author/publisher copy, redistribution terms requiring review.
\end{itemize}
\end{document}
